% GENERATED FILE. Edit canonical lessons, then run npm run generate:books.
% canonical-source-hash: fnv1a64:9c455c38cbd94f00
% canonical-lessons: SA-C152-kalah, SA-C152-varah, SA-C152-tithih, SA-C152-dinankah, SA-C152-satabdi

\chapter{Period, Turn, Lunar day, Calendar date, Century}
\label{ch:sa-period-turn-lunar-day-calendar-date-century}
\hlchaptermodality{152}

\begin{chapteropening}
\textbf{By the end of this chapter:} \emph{I can say a period, a turn, a lunar day, a calendar date and a century.}

Complete the last lesson of chapter 152: I can say a period, a turn, a lunar day, a calendar date and a century.
\end{chapteropening}

\section[\texorpdfstring{\sk{कालः} (kālaḥ)}{कालः (kālaḥ)}]{\sk{कालः} (kālaḥ) — a period, time}

\label{lesson:SA-C152-kalah}



\begin{quote}
Before the new one: say the Sanskrit for jealousy, then the Sanskrit for pride.
\end{quote}

\subsection*{You'll want to know: \sk{कालः}}
\textbf{\sk{कालः}} — \emph{kālaḥ} — \textquotedblleft{}a period, time\textquotedblright{}.

\begin{grammarlens}[title={What you've built}]
One more word for this chapter.
\end{grammarlens}

\subsection*{Your turn}
\begin{itemize}
\raggedright
  \item \emph{Say it:} \emph{kālaḥ}
  \item \emph{Say it:} \emph{kālaḥ}, once more
  \item \emph{Say it:} say \emph{kālaḥ}
  \item \emph{Recall:} say the Sanskrit for jealousy, then the Sanskrit for pride, then say \emph{kālaḥ} again
\end{itemize}

\begin{tcolorbox}[breakable,colback=teal!4,colframe=teal!35!black,title={Before you move on}]
What does \sk{कालः} mean? (\textquotedblleft{}A period, time\textquotedblright{}.) Say it once more.
\end{tcolorbox}
\section[\texorpdfstring{\sk{वारः} (vāraḥ)}{वारः (vāraḥ)}]{\sk{वारः} (vāraḥ) — a turn, a day of the week}

\label{lesson:SA-C152-varah}



\begin{quote}
Before the new one: say the Sanskrit for pride, then the Sanskrit for a period.
\end{quote}

\subsection*{You'll want to know: \sk{वारः}}
\textbf{\sk{वारः}} — \emph{vāraḥ} — \textquotedblleft{}a turn, a day of the week\textquotedblright{}.

\begin{grammarlens}[title={What you've built}]
One more word for this chapter.
\end{grammarlens}

\subsection*{Your turn}
\begin{itemize}
\raggedright
  \item \emph{Say it:} \emph{vāraḥ}
  \item \emph{Say it:} \emph{vāraḥ}, once more
  \item \emph{Say it:} say \emph{vāraḥ}
  \item \emph{Recall:} say the Sanskrit for pride, then the Sanskrit for a period, then say \emph{vāraḥ} again
\end{itemize}

\begin{tcolorbox}[breakable,colback=teal!4,colframe=teal!35!black,title={Before you move on}]
What does \sk{वारः} mean? (\textquotedblleft{}A turn, a day of the week\textquotedblright{}.) Say it once more.
\end{tcolorbox}
\section[\texorpdfstring{\sk{तिथिः} (tithiḥ)}{तिथिः (tithiḥ)}]{\sk{तिथिः} (tithiḥ) — a lunar day, a date}

\label{lesson:SA-C152-tithih}



\begin{quote}
Before the new one: say the Sanskrit for a period, then the Sanskrit for a turn.
\end{quote}

\subsection*{You'll want to know: \sk{तिथिः}}
\textbf{\sk{तिथिः}} — \emph{tithiḥ} — \textquotedblleft{}a lunar day, a date\textquotedblright{}.

\begin{grammarlens}[title={What you've built}]
One more word for this chapter.
\end{grammarlens}

\subsection*{Your turn}
\begin{itemize}
\raggedright
  \item \emph{Say it:} \emph{tithiḥ}
  \item \emph{Say it:} \emph{tithiḥ}, once more
  \item \emph{Say it:} say \emph{tithiḥ}
  \item \emph{Recall:} say the Sanskrit for a period, then the Sanskrit for a turn, then say \emph{tithiḥ} again
\end{itemize}

\begin{tcolorbox}[breakable,colback=teal!4,colframe=teal!35!black,title={Before you move on}]
What does \sk{तिथिः} mean? (\textquotedblleft{}A lunar day, a date\textquotedblright{}.) Say it once more.
\end{tcolorbox}
\section[\texorpdfstring{\sk{दिनांकः} (dināṅkaḥ)}{दिनांकः (dināṅkaḥ)}]{\sk{दिनांकः} (dināṅkaḥ) — a calendar date}

\label{lesson:SA-C152-dinankah}



\begin{quote}
Before the new one: say the Sanskrit for a turn, then the Sanskrit for a lunar day.
\end{quote}

\subsection*{You'll want to know: \sk{दिनांकः}}
\textbf{\sk{दिनांकः}} — \emph{dināṅkaḥ} — \textquotedblleft{}a calendar date\textquotedblright{}.

\begin{grammarlens}[title={What you've built}]
One more word for this chapter.
\end{grammarlens}

\subsection*{Your turn}
\begin{itemize}
\raggedright
  \item \emph{Say it:} \emph{dināṅkaḥ}
  \item \emph{Say it:} \emph{dināṅkaḥ}, once more
  \item \emph{Say it:} say \emph{dināṅkaḥ}
  \item \emph{Recall:} say the Sanskrit for a turn, then the Sanskrit for a lunar day, then say \emph{dināṅkaḥ} again
\end{itemize}

\begin{tcolorbox}[breakable,colback=teal!4,colframe=teal!35!black,title={Before you move on}]
What does \sk{दिनांकः} mean? (\textquotedblleft{}A calendar date\textquotedblright{}.) Say it once more.
\end{tcolorbox}
\section[\texorpdfstring{\sk{शताब्दी} (śatābdī)}{शताब्दी (śatābdī)}]{\sk{शताब्दी} (śatābdī) — a century}

\label{lesson:SA-C152-satabdi}



\begin{quote}
Before the new one: say the Sanskrit for a lunar day, then the Sanskrit for a calendar date.
\end{quote}

\subsection*{You'll want to know: \sk{शताब्दी}}
\textbf{\sk{शताब्दी}} — \emph{śatābdī} — \textquotedblleft{}a century\textquotedblright{}.

\begin{grammarlens}[title={What you've built}]
One more word for this chapter.
\end{grammarlens}

\subsection*{Your turn}
\begin{itemize}
\raggedright
  \item \emph{Say it:} \emph{śatābdī}
  \item \emph{Say it:} \emph{śatābdī}, once more
  \item \emph{Say it:} say \emph{śatābdī}
  \item \emph{Recall:} say the Sanskrit for a lunar day, then the Sanskrit for a calendar date, then say \emph{śatābdī} again
\end{itemize}

\begin{tcolorbox}[breakable,colback=teal!4,colframe=teal!35!black,title={Before you move on}]
What does \sk{शताब्दी} mean? (\textquotedblleft{}A century\textquotedblright{}.) Say it once more.
\end{tcolorbox}
